
\documentclass{article}
\usepackage{colortbl}
\usepackage{makecell}
\usepackage{multirow}
\usepackage{supertabular}

\begin{document}

\newcounter{utterance}

\centering \large Interaction Transcript for game `clean\_up', experiment `2\_hard\_5obj\_fr', episode 0 with gemma{-}4{-}26b{-}a4b{-}it.
\vspace{24pt}

{ \footnotesize  \setcounter{utterance}{1}
\setlength{\tabcolsep}{0pt}
\begin{supertabular}{c@{$\;$}|p{.15\linewidth}@{}p{.15\linewidth}p{.15\linewidth}p{.15\linewidth}p{.15\linewidth}p{.15\linewidth}}
    \# & $\;$A & \multicolumn{4}{c}{Game Master} & $\;\:$B\\
    \hline

    \theutterance \stepcounter{utterance}  
    & & \multicolumn{4}{p{0.6\linewidth}}{
        \cellcolor[rgb]{0.9,0.9,0.9}{
            \makecell[{{p{\linewidth}}}]{
                \texttt{\tiny{[P1$\langle$GM]}}
                \texttt{Je suis votre maître de jeu, et vous jouez à un jeu collaboratif avec la grille suivante comme plateau de jeu :} \\
\\ 
\texttt{```} \\
\texttt{ 1234567} \\
\texttt{╔══╤════╗} \\
\texttt{║◌◌│◌◌◌H║ 1} \\
\texttt{╟──┤◌◌┌─╢ 2} \\
\texttt{║CI├──┤◌║ 3} \\
\texttt{╟──┘◌◌├─╢ 4} \\
\texttt{║◌◌◌◌◌│◌║ 5} \\
\texttt{╟─┐W┌─┼─╢ 6} \\
\texttt{║◌│T│◌│◌║ 7} \\
\texttt{╚═╧═╧═╧═╝} \\
\\ 
\texttt{```} \\
\\ 
\texttt{* Le bord supérieur affiche des coordonnées x croissantes vers la droite, et le bord droit des coordonnées y croissantes vers le bas.} \\
\texttt{* Les objets suivants sont placés aléatoirement sur votre grille : 'W', 'I', 'T', 'C', 'H'.} \\
\\ 
\\ 
\texttt{L'autre joueur voit une variante du plateau de jeu, où les objets sont placés à des emplacements aléatoires différents. Vous ne pouvez pas voir le plateau de l'autre joueur et il ne peut pas voir le vôtre.} \\
\\ 
\\ 
\texttt{**Objectif :**} \\
\\ 
\texttt{Les deux joueurs doivent déplacer les objets sur leur arrière{-}plan respectif de manière à ce que des objets identiques se retrouvent aux mêmes coordonnées. Vous devez communiquer avec l'autre joueur pour convenir d'une configuration d'objectif commune.} \\
\\ 
\\ 
\texttt{**Règles :**} \\
\\ 
\texttt{* À chaque tour, vous pouvez envoyer exactement l'une des deux commandes suivantes :} \\
\texttt{ 1. `DIRE: <MESSAGE>` : pour envoyer un message (tout ce qui se trouve jusqu'au prochain saut de ligne) à l'autre joueur. Je le transmettrai à votre partenaire.} \\
\texttt{2. `DÉPLACER: <OBJET>, (<X>, <Y>)` : pour déplacer un objet vers une nouvelle position, où `<X>` est la colonne et `<Y>` est la ligne. Je vous informerai si votre déplacement a réussi ou non.} \\
\texttt{* Si vous ne respectez pas le format, ou si vous envoyez plusieurs commandes à la fois, je dois vous pénaliser.} \\
\texttt{* Si les deux joueurs accumulent plus de 12 pénalités, vous perdez tous les deux la partie.} \\
\texttt{* Il est essentiel que vous communiquiez avec l'autre joueur concernant votre état d'objectif ! La *seule* façon de transmettre votre stratégie à l'autre joueur est d'utiliser la commande `DIRE: <MESSAGE>` !} \\
\\ 
\\ 
\texttt{**Déplacement d'objets**} \\
\\ 
\texttt{* Vous ne pouvez déplacer des objets que vers des cellules situées à l'intérieur des limites de la grille. La cellule cible doit être vide, c'est{-}à{-}dire qu'elle ne doit contenir que le symbole '◌'.} \\
\texttt{* Si vous essayez de déplacer un objet vers un endroit qui n'est pas vide, ou si vous essayez de le déplacer en dehors de la grille, je dois vous pénaliser. Vous avez droit à un nouvel essai.} \\
\texttt{* Avant d'effectuer un déplacement, vérifiez bien que l'emplacement cible est vide et qu'il ne contient aucune lettre, aucun cadre ou aucune ligne !} \\
\\ 
\\ 
\texttt{**Fin de partie**} \\
\\ 
\texttt{Si vous pensez avoir atteint l'objectif d'aligner tous les objets, vous pouvez demander à l'autre joueur de terminer la partie en envoyant `DIRE: terminé ?`. Si l'autre joueur vous demande de terminer la partie et que vous répondez `DIRE: terminé !`, la partie se termine.} \\
\\ 
\texttt{Les deux joueurs gagnent si la partie se termine en 20 tours, un tour étant défini comme l'envoi d'une commande valide par chacun des deux joueurs.} \\
\\ 
\texttt{**Score :**} \\
\\ 
\texttt{Plus les objets identiques sont proches les uns des autres sur les deux plateaux de jeu, plus vous obtenez de points. Les pénalités réduisent vos points. Pouvez{-}vous battre le record ?} \\
\\ 
\\ 
\texttt{Envoyez un message à l'autre joueur pour commencer le jeu !} \\
            }
        }
    }
    & \\ \\

    \theutterance \stepcounter{utterance}  
    & \multicolumn{4}{p{0.6\linewidth}}{
        \cellcolor[rgb]{0.9,0.9,0.9}{
            \makecell[{{p{\linewidth}}}]{
                \texttt{\tiny{[P1$\rangle$GM]}}
                \texttt{DIRE: Bonjour ! Je vois les objets suivants sur mon plateau : C(1,3), I(2,3), H(7,1), W(3,6) et T(3,7). Pour gagner, nous devons aligner nos objets aux mêmes coordonnées. Proposes{-}tu une configuration cible (par exemple, une ligne ou une colonne spécifique) pour que nous puissions commencer à déplacer nos objets ?} \\
            }
        }
    }
    & & \\ \\

    \theutterance \stepcounter{utterance}  
    & & & \multicolumn{2}{p{0.3\linewidth}}{
        \cellcolor[rgb]{0.9,0.9,0.9}{
            \makecell[{{p{\linewidth}}}]{
                \texttt{\tiny{[GM$|$GM]}}
                \texttt{DIRE: Bonjour ! Je vois les objets suivants sur mon plateau : C(1,3), I(2,3), H(7,1), W(3,6) et T(3,7). Pour gagner, nous devons aligner nos objets aux mêmes coordonnées. Proposes{-}tu une configuration cible (par exemple, une ligne ou une colonne spécifique) pour que nous puissions commencer à déplacer nos objets ?} \\
            }
        }
    }
    & & \\ \\

    \theutterance \stepcounter{utterance}  
    & & \multicolumn{4}{p{0.6\linewidth}}{
        \cellcolor[rgb]{0.9,0.9,0.9}{
            \makecell[{{p{\linewidth}}}]{
                \texttt{\tiny{[GM$\rangle$P2]}}
                \texttt{Je suis votre maître de jeu, et vous jouez à un jeu collaboratif avec la grille suivante comme plateau de jeu :} \\
\\ 
\texttt{```} \\
\texttt{ 1234567} \\
\texttt{╔══╤════╗} \\
\texttt{║◌W│◌◌◌◌║ 1} \\
\texttt{╟──┤◌H┌─╢ 2} \\
\texttt{║◌◌├──┤◌║ 3} \\
\texttt{╟──┘◌◌├─╢ 4} \\
\texttt{║◌◌◌◌◌│◌║ 5} \\
\texttt{╟─┐C┌─┼─╢ 6} \\
\texttt{║I│◌│T│◌║ 7} \\
\texttt{╚═╧═╧═╧═╝} \\
\\ 
\texttt{```} \\
\\ 
\texttt{* Le bord supérieur affiche des coordonnées x croissantes vers la droite, et le bord droit des coordonnées y croissantes vers le bas.} \\
\texttt{* Les objets suivants sont placés aléatoirement sur votre grille : 'W', 'I', 'T', 'C', 'H'.} \\
\\ 
\\ 
\texttt{L'autre joueur voit une variante du plateau de jeu, où les objets sont placés à des emplacements aléatoires différents. Vous ne pouvez pas voir le plateau de l'autre joueur et il ne peut pas voir le vôtre.} \\
\\ 
\\ 
\texttt{**Objectif :**} \\
\\ 
\texttt{Les deux joueurs doivent déplacer les objets sur leur arrière{-}plan respectif de manière à ce que des objets identiques se retrouvent aux mêmes coordonnées. Vous devez communiquer avec l'autre joueur pour convenir d'une configuration d'objectif commune.} \\
\\ 
\\ 
\texttt{**Règles :**} \\
\\ 
\texttt{* À chaque tour, vous pouvez envoyer exactement l'une des deux commandes suivantes :} \\
\texttt{ 1. `DIRE: <MESSAGE>` : pour envoyer un message (tout ce qui se trouve jusqu'au prochain saut de ligne) à l'autre joueur. Je le transmettrai à votre partenaire.} \\
\texttt{2. `DÉPLACER: <OBJET>, (<X>, <Y>)` : pour déplacer un objet vers une nouvelle position, où `<X>` est la colonne et `<Y>` est la ligne. Je vous informerai si votre déplacement a réussi ou non.} \\
\texttt{* Si vous ne respectez pas le format, ou si vous envoyez plusieurs commandes à la fois, je dois vous pénaliser.} \\
\texttt{* Si les deux joueurs accumulent plus de 12 pénalités, vous perdez tous les deux la partie.} \\
\texttt{* Il est essentiel que vous communiquiez avec l'autre joueur concernant votre état d'objectif ! La *seule* façon de transmettre votre stratégie à l'autre joueur est d'utiliser la commande `DIRE: <MESSAGE>` !} \\
\\ 
\\ 
\texttt{**Déplacement d'objets**} \\
\\ 
\texttt{* Vous ne pouvez déplacer des objets que vers des cellules situées à l'intérieur des limites de la grille. La cellule cible doit être vide, c'est{-}à{-}dire qu'elle ne doit contenir que le symbole '◌'.} \\
\texttt{* Si vous essayez de déplacer un objet vers un endroit qui n'est pas vide, ou si vous essayez de le déplacer en dehors de la grille, je dois vous pénaliser. Vous avez droit à un nouvel essai.} \\
\texttt{* Avant d'effectuer un déplacement, vérifiez bien que l'emplacement cible est vide et qu'il ne contient aucune lettre, aucun cadre ou aucune ligne !} \\
\\ 
\\ 
\texttt{**Fin de partie**} \\
\\ 
\texttt{Si vous pensez avoir atteint l'objectif d'aligner tous les objets, vous pouvez demander à l'autre joueur de terminer la partie en envoyant `DIRE: terminé ?`. Si l'autre joueur vous demande de terminer la partie et que vous répondez `DIRE: terminé !`, la partie se termine.} \\
\\ 
\texttt{Les deux joueurs gagnent si la partie se termine en 20 tours, un tour étant défini comme l'envoi d'une commande valide par chacun des deux joueurs.} \\
\\ 
\texttt{**Score :**} \\
\\ 
\texttt{Plus les objets identiques sont proches les uns des autres sur les deux plateaux de jeu, plus vous obtenez de points. Les pénalités réduisent vos points. Pouvez{-}vous battre le record ?} \\
\\ 
\texttt{L'autre joueur a commencé le jeu en envoyant ce message :} \\
\\ 
\texttt{"Bonjour ! Je vois les objets suivants sur mon plateau : C(1,3), I(2,3), H(7,1), W(3,6) et T(3,7). Pour gagner, nous devons aligner nos objets aux mêmes coordonnées. Proposes{-}tu une configuration cible (par exemple, une ligne ou une colonne spécifique) pour que nous puissions commencer à déplacer nos objets ?"} \\
\\ 
\texttt{Quelle est votre première commande ?} \\
            }
        }
    }
    & \\ \\

    \theutterance \stepcounter{utterance}  
    & & & \multicolumn{4}{p{0.6\linewidth}}{
        \cellcolor[rgb]{0.9,0.9,0.9}{
            \makecell[{{p{\linewidth}}}]{
                \texttt{\tiny{[GM$\langle$P2]}}
                \texttt{DIRE: Bonjour ! Merci pour ces informations. Sur mon plateau, j'ai : W(2,1), H(3,2), C(3,6), I(1,7) et T(4,7). Pour simplifier les déplacements, je propose que nous essayions d'aligner tous les objets sur la ligne 7. Par exemple : I(1,7), C(2,7), T(3,7), W(4,7) et H(5,7). Qu'en penses{-}tu ?} \\
            }
        }
    }
     \\ \\

    \theutterance \stepcounter{utterance}  
    & & & \multicolumn{2}{p{0.3\linewidth}}{
        \cellcolor[rgb]{0.9,0.9,0.9}{
            \makecell[{{p{\linewidth}}}]{
                \texttt{\tiny{[GM$|$GM]}}
                \texttt{DIRE: Bonjour ! Merci pour ces informations. Sur mon plateau, j'ai : W(2,1), H(3,2), C(3,6), I(1,7) et T(4,7). Pour simplifier les déplacements, je propose que nous essayions d'aligner tous les objets sur la ligne 7. Par exemple : I(1,7), C(2,7), T(3,7), W(4,7) et H(5,7). Qu'en penses{-}tu ?} \\
            }
        }
    }
    & & \\ \\

    \theutterance \stepcounter{utterance}  
    & & \multicolumn{4}{p{0.6\linewidth}}{
        \cellcolor[rgb]{0.9,0.9,0.9}{
            \makecell[{{p{\linewidth}}}]{
                \texttt{\tiny{[P1$\langle$GM]}}
                \texttt{Votre message a été transmis à l'autre joueur.} \\
\\ 
\\ 
\texttt{Vous jouez actuellement le tour 1 sur un maximum de 20.} \\
\texttt{Vous avez collectivement accumulé 0 pénalités sur 12.} \\
\texttt{L'autre joueur a envoyé ce message :} \\
\\ 
\\ 
\texttt{"Bonjour ! Merci pour ces informations. Sur mon plateau, j'ai : W(2,1), H(3,2), C(3,6), I(1,7) et T(4,7). Pour simplifier les déplacements, je propose que nous essayions d'aligner tous les objets sur la ligne 7. Par exemple : I(1,7), C(2,7), T(3,7), W(4,7) et H(5,7). Qu'en penses{-}tu ?"} \\
\\ 
\\ 
\texttt{Quelle est votre prochaine commande ?} \\
            }
        }
    }
    & \\ \\

    \theutterance \stepcounter{utterance}  
    & \multicolumn{4}{p{0.6\linewidth}}{
        \cellcolor[rgb]{0.9,0.9,0.9}{
            \makecell[{{p{\linewidth}}}]{
                \texttt{\tiny{[P1$\rangle$GM]}}
                \texttt{DIRE: C'est une excellente idée ! La ligne 7 est très accessible pour moi. Cependant, je remarque que sur mon plateau, T est déjà en (3,7). Si nous suivons ta proposition, T serait en (3,7) pour nous deux, ce qui est parfait.} \\
\\ 
\texttt{Voici le plan pour moi :} \\
\texttt{I(1,7), C(2,7), T(3,7), W(4,7), H(5,7).} \\
\\ 
\texttt{Je vais commencer par déplacer mon objet 'I' vers (1,7) si la case est libre. On commence ?} \\
            }
        }
    }
    & & \\ \\

    \theutterance \stepcounter{utterance}  
    & & & \multicolumn{2}{p{0.3\linewidth}}{
        \cellcolor[rgb]{0.9,0.9,0.9}{
            \makecell[{{p{\linewidth}}}]{
                \texttt{\tiny{[GM$|$GM]}}
                \texttt{DIRE: C'est une excellente idée ! La ligne 7 est très accessible pour moi. Cependant, je remarque que sur mon plateau, T est déjà en (3,7). Si nous suivons ta proposition, T serait en (3,7) pour nous deux, ce qui est parfait.} \\
\\ 
\texttt{Voici le plan pour moi :} \\
\texttt{I(1,7), C(2,7), T(3,7), W(4,7), H(5,7).} \\
\\ 
\texttt{Je vais commencer par déplacer mon objet 'I' vers (1,7) si la case est libre. On commence ?} \\
            }
        }
    }
    & & \\ \\

    \theutterance \stepcounter{utterance}  
    & & & \multicolumn{2}{p{0.3\linewidth}}{
        \cellcolor[rgb]{0.9,0.9,0.9}{
            \makecell[{{p{\linewidth}}}]{
                \texttt{\tiny{[GM$|$GM]}}
                \texttt{Invalid format: tail is not empty: '} \\
\\ 
\texttt{Voici le plan pour moi :} \\
\texttt{I(1,7), C(2,7), T(3,7), W(4,7), H(5,7).} \\
\\ 
\texttt{Je vais commencer par déplacer mon objet 'I' vers (1,7) si la case est libre. On commence ?} \\
            }
        }
    }
    & & \\ \\

    \theutterance \stepcounter{utterance}  
    & & \multicolumn{4}{p{0.6\linewidth}}{
        \cellcolor[rgb]{0.9,0.9,0.9}{
            \makecell[{{p{\linewidth}}}]{
                \texttt{\tiny{[P1$\langle$GM]}}
                \texttt{Pénalité : Votre message ne doit rien contenir après la commande !} \\
\\ 
\texttt{Assurez{-}vous que votre réponse ne contienne que `DIRE: <MESSAGE>` ou `DÉPLACER: <OBJET>, (<X>, <Y>)`, et rien d'autre !} \\
\\ 
\texttt{Vous avez collectivement accumulé 1 pénalités sur 12.} \\
\texttt{Veuillez réessayer !} \\
            }
        }
    }
    & \\ \\

    \theutterance \stepcounter{utterance}  
    & \multicolumn{4}{p{0.6\linewidth}}{
        \cellcolor[rgb]{0.9,0.9,0.9}{
            \makecell[{{p{\linewidth}}}]{
                \texttt{\tiny{[P1$\rangle$GM]}}
                \texttt{DIRE: Désolé pour l'erreur. Je suis d'accord pour la ligne 7 : I(1,7), C(2,7), T(3,7), W(4,7), H(5,7). Je commence par déplacer mon I vers (1,7).} \\
            }
        }
    }
    & & \\ \\

    \theutterance \stepcounter{utterance}  
    & & & \multicolumn{2}{p{0.3\linewidth}}{
        \cellcolor[rgb]{0.9,0.9,0.9}{
            \makecell[{{p{\linewidth}}}]{
                \texttt{\tiny{[GM$|$GM]}}
                \texttt{DIRE: Désolé pour l'erreur. Je suis d'accord pour la ligne 7 : I(1,7), C(2,7), T(3,7), W(4,7), H(5,7). Je commence par déplacer mon I vers (1,7).} \\
            }
        }
    }
    & & \\ \\

    \theutterance \stepcounter{utterance}  
    & & \multicolumn{4}{p{0.6\linewidth}}{
        \cellcolor[rgb]{0.9,0.9,0.9}{
            \makecell[{{p{\linewidth}}}]{
                \texttt{\tiny{[GM$\rangle$P2]}}
                \texttt{Votre message a été transmis à l'autre joueur.} \\
\\ 
\\ 
\texttt{Vous jouez actuellement le tour 2 sur un maximum de 20.} \\
\texttt{Vous avez collectivement accumulé 1 pénalités sur 12.} \\
\texttt{L'autre joueur a envoyé ce message :} \\
\\ 
\\ 
\texttt{"Désolé pour l'erreur. Je suis d'accord pour la ligne 7 : I(1,7), C(2,7), T(3,7), W(4,7), H(5,7). Je commence par déplacer mon I vers (1,7)."} \\
\\ 
\\ 
\texttt{Quelle est votre prochaine commande ?} \\
            }
        }
    }
    & \\ \\

    \theutterance \stepcounter{utterance}  
    & & & \multicolumn{4}{p{0.6\linewidth}}{
        \cellcolor[rgb]{0.9,0.9,0.9}{
            \makecell[{{p{\linewidth}}}]{
                \texttt{\tiny{[GM$\langle$P2]}}
                \texttt{DÉPLACER: I, (1, 7)} \\
            }
        }
    }
     \\ \\

    \theutterance \stepcounter{utterance}  
    & & & \multicolumn{2}{p{0.3\linewidth}}{
        \cellcolor[rgb]{0.9,0.9,0.9}{
            \makecell[{{p{\linewidth}}}]{
                \texttt{\tiny{[GM$|$GM]}}
                \texttt{DÉPLACER: I, (1, 7)} \\
            }
        }
    }
    & & \\ \\

    \theutterance \stepcounter{utterance}  
    & & & \multicolumn{2}{p{0.3\linewidth}}{
        \cellcolor[rgb]{0.9,0.9,0.9}{
            \makecell[{{p{\linewidth}}}]{
                \texttt{\tiny{[GM$|$GM]}}
                \texttt{Pénalité : (1,7) n'est pas vide, mais contient 'I'.} \\
\\ 
\texttt{Vous avez collectivement accumulé 2 pénalités sur 12.} \\
\texttt{Veuillez réessayer !} \\
            }
        }
    }
    & & \\ \\

    \theutterance \stepcounter{utterance}  
    & & \multicolumn{4}{p{0.6\linewidth}}{
        \cellcolor[rgb]{0.9,0.9,0.9}{
            \makecell[{{p{\linewidth}}}]{
                \texttt{\tiny{[GM$\rangle$P2]}}
                \texttt{Pénalité : (1,7) n'est pas vide, mais contient 'I'.} \\
\\ 
\texttt{Vous avez collectivement accumulé 2 pénalités sur 12.} \\
\texttt{Veuillez réessayer !} \\
            }
        }
    }
    & \\ \\

    \theutterance \stepcounter{utterance}  
    & & & \multicolumn{4}{p{0.6\linewidth}}{
        \cellcolor[rgb]{0.9,0.9,0.9}{
            \makecell[{{p{\linewidth}}}]{
                \texttt{\tiny{[GM$\langle$P2]}}
                \texttt{DÉPLACER: I, (1, 7) est déjà là, je voulais dire que je vais déplacer un autre objet. Je vais commencer par déplacer mon T vers la position cible (3,7).} \\
\\ 
\texttt{DÉPLACER: T, (3, 7)} \\
            }
        }
    }
     \\ \\

    \theutterance \stepcounter{utterance}  
    & & & \multicolumn{2}{p{0.3\linewidth}}{
        \cellcolor[rgb]{0.9,0.9,0.9}{
            \makecell[{{p{\linewidth}}}]{
                \texttt{\tiny{[GM$|$GM]}}
                \texttt{DÉPLACER: I, (1, 7) est déjà là, je voulais dire que je vais déplacer un autre objet. Je vais commencer par déplacer mon T vers la position cible (3,7).} \\
\\ 
\texttt{DÉPLACER: T, (3, 7)} \\
            }
        }
    }
    & & \\ \\

    \theutterance \stepcounter{utterance}  
    & & & \multicolumn{2}{p{0.3\linewidth}}{
        \cellcolor[rgb]{0.9,0.9,0.9}{
            \makecell[{{p{\linewidth}}}]{
                \texttt{\tiny{[GM$|$GM]}}
                \texttt{Invalid format: head is not empty: 'DÉPLACER: I, (1, 7) est déjà là, je voulais dire que je vais déplacer un autre objet. Je vais commencer par déplacer mon T vers la position cible (3,7).} \\
\\ 
\texttt{'} \\
            }
        }
    }
    & & \\ \\

    \theutterance \stepcounter{utterance}  
    & & \multicolumn{4}{p{0.6\linewidth}}{
        \cellcolor[rgb]{0.9,0.9,0.9}{
            \makecell[{{p{\linewidth}}}]{
                \texttt{\tiny{[GM$\rangle$P2]}}
                \texttt{Pénalité : Votre message ne doit rien contenir avant la commande !} \\
\\ 
\texttt{Assurez{-}vous que votre réponse ne contienne que `DIRE: <MESSAGE>` ou `DÉPLACER: <OBJET>, (<X>, <Y>)`, et rien d'autre !} \\
\\ 
\texttt{Vous avez collectivement accumulé 3 pénalités sur 12.} \\
\texttt{Veuillez réessayer !} \\
            }
        }
    }
    & \\ \\

    \theutterance \stepcounter{utterance}  
    & & & \multicolumn{4}{p{0.6\linewidth}}{
        \cellcolor[rgb]{0.9,0.9,0.9}{
            \makecell[{{p{\linewidth}}}]{
                \texttt{\tiny{[GM$\langle$P2]}}
                \texttt{DÉPLACER: T, (3, 7)} \\
            }
        }
    }
     \\ \\

    \theutterance \stepcounter{utterance}  
    & & & \multicolumn{2}{p{0.3\linewidth}}{
        \cellcolor[rgb]{0.9,0.9,0.9}{
            \makecell[{{p{\linewidth}}}]{
                \texttt{\tiny{[GM$|$GM]}}
                \texttt{DÉPLACER: T, (3, 7)} \\
            }
        }
    }
    & & \\ \\

    \theutterance \stepcounter{utterance}  
    & & & \multicolumn{2}{p{0.3\linewidth}}{
        \cellcolor[rgb]{0.9,0.9,0.9}{
            \makecell[{{p{\linewidth}}}]{
                \texttt{\tiny{[GM$|$GM]}}
                \texttt{T moved to (3, 7)} \\
            }
        }
    }
    & & \\ \\

    \theutterance \stepcounter{utterance}  
    & & \multicolumn{4}{p{0.6\linewidth}}{
        \cellcolor[rgb]{0.9,0.9,0.9}{
            \makecell[{{p{\linewidth}}}]{
                \texttt{\tiny{[P1$\langle$GM]}}
                \texttt{Votre message a été transmis à l'autre joueur.} \\
\\ 
\\ 
\texttt{Vous jouez actuellement le tour 2 sur un maximum de 20.} \\
\texttt{Vous avez collectivement accumulé 3 pénalités sur 12.} \\
\texttt{L'autre joueur a déplacé un objet sur sa grille.} \\
\\ 
\\ 
\texttt{Quelle est votre prochaine commande ?} \\
            }
        }
    }
    & \\ \\

    \theutterance \stepcounter{utterance}  
    & \multicolumn{4}{p{0.6\linewidth}}{
        \cellcolor[rgb]{0.9,0.9,0.9}{
            \makecell[{{p{\linewidth}}}]{
                \texttt{\tiny{[P1$\rangle$GM]}}
                \texttt{DÉPLACER: I, (1, 7)} \\
            }
        }
    }
    & & \\ \\

    \theutterance \stepcounter{utterance}  
    & & & \multicolumn{2}{p{0.3\linewidth}}{
        \cellcolor[rgb]{0.9,0.9,0.9}{
            \makecell[{{p{\linewidth}}}]{
                \texttt{\tiny{[GM$|$GM]}}
                \texttt{DÉPLACER: I, (1, 7)} \\
            }
        }
    }
    & & \\ \\

    \theutterance \stepcounter{utterance}  
    & & & \multicolumn{2}{p{0.3\linewidth}}{
        \cellcolor[rgb]{0.9,0.9,0.9}{
            \makecell[{{p{\linewidth}}}]{
                \texttt{\tiny{[GM$|$GM]}}
                \texttt{I moved to (1, 7)} \\
            }
        }
    }
    & & \\ \\

    \theutterance \stepcounter{utterance}  
    & & \multicolumn{4}{p{0.6\linewidth}}{
        \cellcolor[rgb]{0.9,0.9,0.9}{
            \makecell[{{p{\linewidth}}}]{
                \texttt{\tiny{[GM$\rangle$P2]}}
                \texttt{Déplacement de 'T' vers (3,7) réussi. Votre grille mise à jour ressemble à ceci :} \\
\texttt{```} \\
\texttt{ 1234567} \\
\texttt{╔══╤════╗} \\
\texttt{║◌W│◌◌◌◌║ 1} \\
\texttt{╟──┤◌H┌─╢ 2} \\
\texttt{║◌◌├──┤◌║ 3} \\
\texttt{╟──┘◌◌├─╢ 4} \\
\texttt{║◌◌◌◌◌│◌║ 5} \\
\texttt{╟─┐C┌─┼─╢ 6} \\
\texttt{║I│T│◌│◌║ 7} \\
\texttt{╚═╧═╧═╧═╝} \\
\texttt{```} \\
\texttt{Vous jouez actuellement le tour 3 sur un maximum de 20.} \\
\texttt{Vous avez collectivement accumulé 3 pénalités sur 12.} \\
\texttt{L'autre joueur a déplacé un objet sur sa grille.} \\
\\ 
\\ 
\texttt{Quelle est votre prochaine commande ?} \\
            }
        }
    }
    & \\ \\

    \theutterance \stepcounter{utterance}  
    & & & \multicolumn{4}{p{0.6\linewidth}}{
        \cellcolor[rgb]{0.9,0.9,0.9}{
            \makecell[{{p{\linewidth}}}]{
                \texttt{\tiny{[GM$\langle$P2]}}
                \texttt{DÉPLACER: C, (2, 7)} \\
            }
        }
    }
     \\ \\

    \theutterance \stepcounter{utterance}  
    & & & \multicolumn{2}{p{0.3\linewidth}}{
        \cellcolor[rgb]{0.9,0.9,0.9}{
            \makecell[{{p{\linewidth}}}]{
                \texttt{\tiny{[GM$|$GM]}}
                \texttt{DÉPLACER: C, (2, 7)} \\
            }
        }
    }
    & & \\ \\

    \theutterance \stepcounter{utterance}  
    & & & \multicolumn{2}{p{0.3\linewidth}}{
        \cellcolor[rgb]{0.9,0.9,0.9}{
            \makecell[{{p{\linewidth}}}]{
                \texttt{\tiny{[GM$|$GM]}}
                \texttt{Pénalité : (2,7) n'est pas vide, mais contient '│'.} \\
\\ 
\texttt{Vous avez collectivement accumulé 4 pénalités sur 12.} \\
\texttt{Veuillez réessayer !} \\
            }
        }
    }
    & & \\ \\

    \theutterance \stepcounter{utterance}  
    & & \multicolumn{4}{p{0.6\linewidth}}{
        \cellcolor[rgb]{0.9,0.9,0.9}{
            \makecell[{{p{\linewidth}}}]{
                \texttt{\tiny{[GM$\rangle$P2]}}
                \texttt{Pénalité : (2,7) n'est pas vide, mais contient '│'.} \\
\\ 
\texttt{Vous avez collectivement accumulé 4 pénalités sur 12.} \\
\texttt{Veuillez réessayer !} \\
            }
        }
    }
    & \\ \\

    \theutterance \stepcounter{utterance}  
    & & & \multicolumn{4}{p{0.6\linewidth}}{
        \cellcolor[rgb]{0.9,0.9,0.9}{
            \makecell[{{p{\linewidth}}}]{
                \texttt{\tiny{[GM$\langle$P2]}}
                \texttt{DÉPLACER: C, (3, 7) est occupé par T, je vais essayer de déplacer C vers (5, 7) pour libérer de l'espace ou avancer. Attends, je vais d'abord déplacer C vers (2, 6) pour ne pas être bloqué par les bordures. Non, restons sur la ligne 7. Je vais déplacer C vers (5, 7).} \\
\\ 
\texttt{DÉPLACER: C, (5, 7)} \\
            }
        }
    }
     \\ \\

    \theutterance \stepcounter{utterance}  
    & & & \multicolumn{2}{p{0.3\linewidth}}{
        \cellcolor[rgb]{0.9,0.9,0.9}{
            \makecell[{{p{\linewidth}}}]{
                \texttt{\tiny{[GM$|$GM]}}
                \texttt{DÉPLACER: C, (3, 7) est occupé par T, je vais essayer de déplacer C vers (5, 7) pour libérer de l'espace ou avancer. Attends, je vais d'abord déplacer C vers (2, 6) pour ne pas être bloqué par les bordures. Non, restons sur la ligne 7. Je vais déplacer C vers (5, 7).} \\
\\ 
\texttt{DÉPLACER: C, (5, 7)} \\
            }
        }
    }
    & & \\ \\

    \theutterance \stepcounter{utterance}  
    & & & \multicolumn{2}{p{0.3\linewidth}}{
        \cellcolor[rgb]{0.9,0.9,0.9}{
            \makecell[{{p{\linewidth}}}]{
                \texttt{\tiny{[GM$|$GM]}}
                \texttt{Invalid format: head is not empty: 'DÉPLACER: C, (3, 7) est occupé par T, je vais essayer de déplacer C vers (5, 7) pour libérer de l'espace ou avancer. Attends, je vais d'abord déplacer C vers (2, 6) pour ne pas être bloqué par les bordures. Non, restons sur la ligne 7. Je vais déplacer C vers (5, 7).} \\
\\ 
\texttt{'} \\
            }
        }
    }
    & & \\ \\

    \theutterance \stepcounter{utterance}  
    & & \multicolumn{4}{p{0.6\linewidth}}{
        \cellcolor[rgb]{0.9,0.9,0.9}{
            \makecell[{{p{\linewidth}}}]{
                \texttt{\tiny{[GM$\rangle$P2]}}
                \texttt{Pénalité : Votre message ne doit rien contenir avant la commande !} \\
\\ 
\texttt{Assurez{-}vous que votre réponse ne contienne que `DIRE: <MESSAGE>` ou `DÉPLACER: <OBJET>, (<X>, <Y>)`, et rien d'autre !} \\
\\ 
\texttt{Vous avez collectivement accumulé 5 pénalités sur 12.} \\
\texttt{Veuillez réessayer !} \\
            }
        }
    }
    & \\ \\

    \theutterance \stepcounter{utterance}  
    & & & \multicolumn{4}{p{0.6\linewidth}}{
        \cellcolor[rgb]{0.9,0.9,0.9}{
            \makecell[{{p{\linewidth}}}]{
                \texttt{\tiny{[GM$\langle$P2]}}
                \texttt{DÉPLACER: C, (5, 7)} \\
            }
        }
    }
     \\ \\

    \theutterance \stepcounter{utterance}  
    & & & \multicolumn{2}{p{0.3\linewidth}}{
        \cellcolor[rgb]{0.9,0.9,0.9}{
            \makecell[{{p{\linewidth}}}]{
                \texttt{\tiny{[GM$|$GM]}}
                \texttt{DÉPLACER: C, (5, 7)} \\
            }
        }
    }
    & & \\ \\

    \theutterance \stepcounter{utterance}  
    & & & \multicolumn{2}{p{0.3\linewidth}}{
        \cellcolor[rgb]{0.9,0.9,0.9}{
            \makecell[{{p{\linewidth}}}]{
                \texttt{\tiny{[GM$|$GM]}}
                \texttt{C moved to (5, 7)} \\
            }
        }
    }
    & & \\ \\

    \theutterance \stepcounter{utterance}  
    & & \multicolumn{4}{p{0.6\linewidth}}{
        \cellcolor[rgb]{0.9,0.9,0.9}{
            \makecell[{{p{\linewidth}}}]{
                \texttt{\tiny{[P1$\langle$GM]}}
                \texttt{Déplacement de 'I' vers (1,7) réussi. Votre grille mise à jour ressemble à ceci :} \\
\texttt{```} \\
\texttt{ 1234567} \\
\texttt{╔══╤════╗} \\
\texttt{║◌◌│◌◌◌H║ 1} \\
\texttt{╟──┤◌◌┌─╢ 2} \\
\texttt{║C◌├──┤◌║ 3} \\
\texttt{╟──┘◌◌├─╢ 4} \\
\texttt{║◌◌◌◌◌│◌║ 5} \\
\texttt{╟─┐W┌─┼─╢ 6} \\
\texttt{║I│T│◌│◌║ 7} \\
\texttt{╚═╧═╧═╧═╝} \\
\texttt{```} \\
\texttt{Vous jouez actuellement le tour 3 sur un maximum de 20.} \\
\texttt{Vous avez collectivement accumulé 5 pénalités sur 12.} \\
\texttt{L'autre joueur a déplacé un objet sur sa grille.} \\
\\ 
\\ 
\texttt{Quelle est votre prochaine commande ?} \\
            }
        }
    }
    & \\ \\

    \theutterance \stepcounter{utterance}  
    & \multicolumn{4}{p{0.6\linewidth}}{
        \cellcolor[rgb]{0.9,0.9,0.9}{
            \makecell[{{p{\linewidth}}}]{
                \texttt{\tiny{[P1$\rangle$GM]}}
                \texttt{DÉPLACER: C, (2, 7)} \\
            }
        }
    }
    & & \\ \\

    \theutterance \stepcounter{utterance}  
    & & & \multicolumn{2}{p{0.3\linewidth}}{
        \cellcolor[rgb]{0.9,0.9,0.9}{
            \makecell[{{p{\linewidth}}}]{
                \texttt{\tiny{[GM$|$GM]}}
                \texttt{DÉPLACER: C, (2, 7)} \\
            }
        }
    }
    & & \\ \\

    \theutterance \stepcounter{utterance}  
    & & & \multicolumn{2}{p{0.3\linewidth}}{
        \cellcolor[rgb]{0.9,0.9,0.9}{
            \makecell[{{p{\linewidth}}}]{
                \texttt{\tiny{[GM$|$GM]}}
                \texttt{Pénalité : (2,7) n'est pas vide, mais contient '│'.} \\
\\ 
\texttt{Vous avez collectivement accumulé 6 pénalités sur 12.} \\
\texttt{Veuillez réessayer !} \\
            }
        }
    }
    & & \\ \\

    \theutterance \stepcounter{utterance}  
    & & \multicolumn{4}{p{0.6\linewidth}}{
        \cellcolor[rgb]{0.9,0.9,0.9}{
            \makecell[{{p{\linewidth}}}]{
                \texttt{\tiny{[P1$\langle$GM]}}
                \texttt{Pénalité : (2,7) n'est pas vide, mais contient '│'.} \\
\\ 
\texttt{Vous avez collectivement accumulé 6 pénalités sur 12.} \\
\texttt{Veuillez réessayer !} \\
            }
        }
    }
    & \\ \\

    \theutterance \stepcounter{utterance}  
    & \multicolumn{4}{p{0.6\linewidth}}{
        \cellcolor[rgb]{0.9,0.9,0.9}{
            \makecell[{{p{\linewidth}}}]{
                \texttt{\tiny{[P1$\rangle$GM]}}
                \texttt{DÉPLACER: C, (2, 3)} \\
            }
        }
    }
    & & \\ \\

    \theutterance \stepcounter{utterance}  
    & & & \multicolumn{2}{p{0.3\linewidth}}{
        \cellcolor[rgb]{0.9,0.9,0.9}{
            \makecell[{{p{\linewidth}}}]{
                \texttt{\tiny{[GM$|$GM]}}
                \texttt{DÉPLACER: C, (2, 3)} \\
            }
        }
    }
    & & \\ \\

    \theutterance \stepcounter{utterance}  
    & & & \multicolumn{2}{p{0.3\linewidth}}{
        \cellcolor[rgb]{0.9,0.9,0.9}{
            \makecell[{{p{\linewidth}}}]{
                \texttt{\tiny{[GM$|$GM]}}
                \texttt{C moved to (2, 3)} \\
            }
        }
    }
    & & \\ \\

    \theutterance \stepcounter{utterance}  
    & & \multicolumn{4}{p{0.6\linewidth}}{
        \cellcolor[rgb]{0.9,0.9,0.9}{
            \makecell[{{p{\linewidth}}}]{
                \texttt{\tiny{[GM$\rangle$P2]}}
                \texttt{Déplacement de 'C' vers (5,7) réussi. Votre grille mise à jour ressemble à ceci :} \\
\texttt{```} \\
\texttt{ 1234567} \\
\texttt{╔══╤════╗} \\
\texttt{║◌W│◌◌◌◌║ 1} \\
\texttt{╟──┤◌H┌─╢ 2} \\
\texttt{║◌◌├──┤◌║ 3} \\
\texttt{╟──┘◌◌├─╢ 4} \\
\texttt{║◌◌◌◌◌│◌║ 5} \\
\texttt{╟─┐◌┌─┼─╢ 6} \\
\texttt{║I│T│C│◌║ 7} \\
\texttt{╚═╧═╧═╧═╝} \\
\texttt{```} \\
\texttt{Vous jouez actuellement le tour 4 sur un maximum de 20.} \\
\texttt{Vous avez collectivement accumulé 6 pénalités sur 12.} \\
\texttt{L'autre joueur a déplacé un objet sur sa grille.} \\
\\ 
\\ 
\texttt{Quelle est votre prochaine commande ?} \\
            }
        }
    }
    & \\ \\

    \theutterance \stepcounter{utterance}  
    & & & \multicolumn{4}{p{0.6\linewidth}}{
        \cellcolor[rgb]{0.9,0.9,0.9}{
            \makecell[{{p{\linewidth}}}]{
                \texttt{\tiny{[GM$\langle$P2]}}
                \texttt{DÉPLACER: W, (4, 7)} \\
            }
        }
    }
     \\ \\

    \theutterance \stepcounter{utterance}  
    & & & \multicolumn{2}{p{0.3\linewidth}}{
        \cellcolor[rgb]{0.9,0.9,0.9}{
            \makecell[{{p{\linewidth}}}]{
                \texttt{\tiny{[GM$|$GM]}}
                \texttt{DÉPLACER: W, (4, 7)} \\
            }
        }
    }
    & & \\ \\

    \theutterance \stepcounter{utterance}  
    & & & \multicolumn{2}{p{0.3\linewidth}}{
        \cellcolor[rgb]{0.9,0.9,0.9}{
            \makecell[{{p{\linewidth}}}]{
                \texttt{\tiny{[GM$|$GM]}}
                \texttt{Pénalité : (4,7) n'est pas vide, mais contient '│'.} \\
\\ 
\texttt{Vous avez collectivement accumulé 7 pénalités sur 12.} \\
\texttt{Veuillez réessayer !} \\
            }
        }
    }
    & & \\ \\

    \theutterance \stepcounter{utterance}  
    & & \multicolumn{4}{p{0.6\linewidth}}{
        \cellcolor[rgb]{0.9,0.9,0.9}{
            \makecell[{{p{\linewidth}}}]{
                \texttt{\tiny{[GM$\rangle$P2]}}
                \texttt{Pénalité : (4,7) n'est pas vide, mais contient '│'.} \\
\\ 
\texttt{Vous avez collectivement accumulé 7 pénalités sur 12.} \\
\texttt{Veuillez réessayer !} \\
            }
        }
    }
    & \\ \\

    \theutterance \stepcounter{utterance}  
    & & & \multicolumn{4}{p{0.6\linewidth}}{
        \cellcolor[rgb]{0.9,0.9,0.9}{
            \makecell[{{p{\linewidth}}}]{
                \texttt{\tiny{[GM$\langle$P2]}}
                \texttt{DÉPLACER: W, (4, 1)} \\
            }
        }
    }
     \\ \\

    \theutterance \stepcounter{utterance}  
    & & & \multicolumn{2}{p{0.3\linewidth}}{
        \cellcolor[rgb]{0.9,0.9,0.9}{
            \makecell[{{p{\linewidth}}}]{
                \texttt{\tiny{[GM$|$GM]}}
                \texttt{DÉPLACER: W, (4, 1)} \\
            }
        }
    }
    & & \\ \\

    \theutterance \stepcounter{utterance}  
    & & & \multicolumn{2}{p{0.3\linewidth}}{
        \cellcolor[rgb]{0.9,0.9,0.9}{
            \makecell[{{p{\linewidth}}}]{
                \texttt{\tiny{[GM$|$GM]}}
                \texttt{W moved to (4, 1)} \\
            }
        }
    }
    & & \\ \\

    \theutterance \stepcounter{utterance}  
    & & \multicolumn{4}{p{0.6\linewidth}}{
        \cellcolor[rgb]{0.9,0.9,0.9}{
            \makecell[{{p{\linewidth}}}]{
                \texttt{\tiny{[P1$\langle$GM]}}
                \texttt{Déplacement de 'C' vers (2,3) réussi. Votre grille mise à jour ressemble à ceci :} \\
\texttt{```} \\
\texttt{ 1234567} \\
\texttt{╔══╤════╗} \\
\texttt{║◌◌│◌◌◌H║ 1} \\
\texttt{╟──┤◌◌┌─╢ 2} \\
\texttt{║◌C├──┤◌║ 3} \\
\texttt{╟──┘◌◌├─╢ 4} \\
\texttt{║◌◌◌◌◌│◌║ 5} \\
\texttt{╟─┐W┌─┼─╢ 6} \\
\texttt{║I│T│◌│◌║ 7} \\
\texttt{╚═╧═╧═╧═╝} \\
\texttt{```} \\
\texttt{Vous jouez actuellement le tour 4 sur un maximum de 20.} \\
\texttt{Vous avez collectivement accumulé 7 pénalités sur 12.} \\
\texttt{L'autre joueur a déplacé un objet sur sa grille.} \\
\\ 
\\ 
\texttt{Quelle est votre prochaine commande ?} \\
            }
        }
    }
    & \\ \\

    \theutterance \stepcounter{utterance}  
    & \multicolumn{4}{p{0.6\linewidth}}{
        \cellcolor[rgb]{0.9,0.9,0.9}{
            \makecell[{{p{\linewidth}}}]{
                \texttt{\tiny{[P1$\rangle$GM]}}
                \texttt{DÉPLACER: C, (2, 7)} \\
            }
        }
    }
    & & \\ \\

    \theutterance \stepcounter{utterance}  
    & & & \multicolumn{2}{p{0.3\linewidth}}{
        \cellcolor[rgb]{0.9,0.9,0.9}{
            \makecell[{{p{\linewidth}}}]{
                \texttt{\tiny{[GM$|$GM]}}
                \texttt{DÉPLACER: C, (2, 7)} \\
            }
        }
    }
    & & \\ \\

    \theutterance \stepcounter{utterance}  
    & & & \multicolumn{2}{p{0.3\linewidth}}{
        \cellcolor[rgb]{0.9,0.9,0.9}{
            \makecell[{{p{\linewidth}}}]{
                \texttt{\tiny{[GM$|$GM]}}
                \texttt{Pénalité : (2,7) n'est pas vide, mais contient '│'.} \\
\\ 
\texttt{Vous avez collectivement accumulé 8 pénalités sur 12.} \\
\texttt{Veuillez réessayer !} \\
            }
        }
    }
    & & \\ \\

    \theutterance \stepcounter{utterance}  
    & & \multicolumn{4}{p{0.6\linewidth}}{
        \cellcolor[rgb]{0.9,0.9,0.9}{
            \makecell[{{p{\linewidth}}}]{
                \texttt{\tiny{[P1$\langle$GM]}}
                \texttt{Pénalité : (2,7) n'est pas vide, mais contient '│'.} \\
\\ 
\texttt{Vous avez collectivement accumulé 8 pénalités sur 12.} \\
\texttt{Veuillez réessayer !} \\
            }
        }
    }
    & \\ \\

    \theutterance \stepcounter{utterance}  
    & \multicolumn{4}{p{0.6\linewidth}}{
        \cellcolor[rgb]{0.9,0.9,0.9}{
            \makecell[{{p{\linewidth}}}]{
                \texttt{\tiny{[P1$\rangle$GM]}}
                \texttt{DÉPLACER: C, (4, 7)} \\
            }
        }
    }
    & & \\ \\

    \theutterance \stepcounter{utterance}  
    & & & \multicolumn{2}{p{0.3\linewidth}}{
        \cellcolor[rgb]{0.9,0.9,0.9}{
            \makecell[{{p{\linewidth}}}]{
                \texttt{\tiny{[GM$|$GM]}}
                \texttt{DÉPLACER: C, (4, 7)} \\
            }
        }
    }
    & & \\ \\

    \theutterance \stepcounter{utterance}  
    & & & \multicolumn{2}{p{0.3\linewidth}}{
        \cellcolor[rgb]{0.9,0.9,0.9}{
            \makecell[{{p{\linewidth}}}]{
                \texttt{\tiny{[GM$|$GM]}}
                \texttt{Pénalité : (4,7) n'est pas vide, mais contient '│'.} \\
\\ 
\texttt{Vous avez collectivement accumulé 9 pénalités sur 12.} \\
\texttt{Veuillez réessayer !} \\
            }
        }
    }
    & & \\ \\

    \theutterance \stepcounter{utterance}  
    & & \multicolumn{4}{p{0.6\linewidth}}{
        \cellcolor[rgb]{0.9,0.9,0.9}{
            \makecell[{{p{\linewidth}}}]{
                \texttt{\tiny{[P1$\langle$GM]}}
                \texttt{Pénalité : (4,7) n'est pas vide, mais contient '│'.} \\
\\ 
\texttt{Vous avez collectivement accumulé 9 pénalités sur 12.} \\
\texttt{Veuillez réessayer !} \\
            }
        }
    }
    & \\ \\

    \theutterance \stepcounter{utterance}  
    & \multicolumn{4}{p{0.6\linewidth}}{
        \cellcolor[rgb]{0.9,0.9,0.9}{
            \makecell[{{p{\linewidth}}}]{
                \texttt{\tiny{[P1$\rangle$GM]}}
                \texttt{DÉPLACER: C, (2, 7)} \\
\\ 
\texttt{Wait, I see the error. The cell (2,7) contains a vertical bar '│'. Looking at the grid:} \\
\texttt{Row 7 is `║I│T│◌│◌║`.} \\
\texttt{The columns are:} \\
\texttt{1: I} \\
\texttt{2: │} \\
\texttt{3: T} \\
\texttt{4: │} \\
\texttt{5: ◌} \\
\texttt{6: │} \\
\texttt{7: ◌} \\
\\ 
\texttt{My target coordinates for the line 7 were: I(1,7), C(2,7), T(3,7), W(4,7), H(5,7).} \\
\texttt{But the grid structure shows that columns 2, 4, and 6 are occupied by the border '│'.} \\
\texttt{I must choose columns that are '◌'.} \\
\texttt{Available '◌' in row 7 are: (5,7) and (7,7).} \\
\\ 
\texttt{Let's re{-}evaluate the target line. If I can only use '◌' cells:} \\
\texttt{Row 7 has '◌' at (5,7) and (7,7).} \\
\texttt{Row 5 is `║◌◌◌◌◌│◌║`. It has '◌' at (1,5), (2,5), (3,5), (4,5), (5,5) and (7,5).} \\
\\ 
\texttt{Let's try to align on Row 5 instead.} \\
\texttt{Target: I(1,5), C(2,5), T(3,5), W(4,5), H(5,5).} \\
\\ 
\texttt{Let's try to move C to (2,5).} \\
\\ 
\texttt{DÉPLACER: C, (2, 5)} \\
            }
        }
    }
    & & \\ \\

    \theutterance \stepcounter{utterance}  
    & & & \multicolumn{2}{p{0.3\linewidth}}{
        \cellcolor[rgb]{0.9,0.9,0.9}{
            \makecell[{{p{\linewidth}}}]{
                \texttt{\tiny{[GM$|$GM]}}
                \texttt{DÉPLACER: C, (2, 7)} \\
\\ 
\texttt{Wait, I see the error. The cell (2,7) contains a vertical bar '│'. Looking at the grid:} \\
\texttt{Row 7 is `║I│T│◌│◌║`.} \\
\texttt{The columns are:} \\
\texttt{1: I} \\
\texttt{2: │} \\
\texttt{3: T} \\
\texttt{4: │} \\
\texttt{5: ◌} \\
\texttt{6: │} \\
\texttt{7: ◌} \\
\\ 
\texttt{My target coordinates for the line 7 were: I(1,7), C(2,7), T(3,7), W(4,7), H(5,7).} \\
\texttt{But the grid structure shows that columns 2, 4, and 6 are occupied by the border '│'.} \\
\texttt{I must choose columns that are '◌'.} \\
\texttt{Available '◌' in row 7 are: (5,7) and (7,7).} \\
\\ 
\texttt{Let's re{-}evaluate the target line. If I can only use '◌' cells:} \\
\texttt{Row 7 has '◌' at (5,7) and (7,7).} \\
\texttt{Row 5 is `║◌◌◌◌◌│◌║`. It has '◌' at (1,5), (2,5), (3,5), (4,5), (5,5) and (7,5).} \\
\\ 
\texttt{Let's try to align on Row 5 instead.} \\
\texttt{Target: I(1,5), C(2,5), T(3,5), W(4,5), H(5,5).} \\
\\ 
\texttt{Let's try to move C to (2,5).} \\
\\ 
\texttt{DÉPLACER: C, (2, 5)} \\
            }
        }
    }
    & & \\ \\

    \theutterance \stepcounter{utterance}  
    & & & \multicolumn{2}{p{0.3\linewidth}}{
        \cellcolor[rgb]{0.9,0.9,0.9}{
            \makecell[{{p{\linewidth}}}]{
                \texttt{\tiny{[GM$|$GM]}}
                \texttt{Invalid format: head is not empty: 'DÉPLACER: C, (2, 7)} \\
\\ 
\texttt{Wait, I see the error. The cell (2,7) contains a vertical bar '│'. Looking at the grid:} \\
\texttt{Row 7 is ║I│T│◌│◌║.} \\
\texttt{The columns are:} \\
\texttt{1: I} \\
\texttt{2: │} \\
\texttt{3: T} \\
\texttt{4: │} \\
\texttt{5: ◌} \\
\texttt{6: │} \\
\texttt{7: ◌} \\
\\ 
\texttt{My target coordinates for the line 7 were: I(1,7), C(2,7), T(3,7), W(4,7), H(5,7).} \\
\texttt{But the grid structure shows that columns 2, 4, and 6 are occupied by the border '│'.} \\
\texttt{I must choose columns that are '◌'.} \\
\texttt{Available '◌' in row 7 are: (5,7) and (7,7).} \\
\\ 
\texttt{Let's re{-}evaluate the target line. If I can only use '◌' cells:} \\
\texttt{Row 7 has '◌' at (5,7) and (7,7).} \\
\texttt{Row 5 is ║◌◌◌◌◌│◌║. It has '◌' at (1,5), (2,5), (3,5), (4,5), (5,5) and (7,5).} \\
\\ 
\texttt{Let's try to align on Row 5 instead.} \\
\texttt{Target: I(1,5), C(2,5), T(3,5), W(4,5), H(5,5).} \\
\\ 
\texttt{Let's try to move C to (2,5).} \\
\\ 
\texttt{'} \\
            }
        }
    }
    & & \\ \\

    \theutterance \stepcounter{utterance}  
    & & \multicolumn{4}{p{0.6\linewidth}}{
        \cellcolor[rgb]{0.9,0.9,0.9}{
            \makecell[{{p{\linewidth}}}]{
                \texttt{\tiny{[P1$\langle$GM]}}
                \texttt{Pénalité : Votre message ne doit rien contenir avant la commande !} \\
\\ 
\texttt{Assurez{-}vous que votre réponse ne contienne que `DIRE: <MESSAGE>` ou `DÉPLACER: <OBJET>, (<X>, <Y>)`, et rien d'autre !} \\
\\ 
\texttt{Vous avez collectivement accumulé 10 pénalités sur 12.} \\
\texttt{Veuillez réessayer !} \\
            }
        }
    }
    & \\ \\

    \theutterance \stepcounter{utterance}  
    & \multicolumn{4}{p{0.6\linewidth}}{
        \cellcolor[rgb]{0.9,0.9,0.9}{
            \makecell[{{p{\linewidth}}}]{
                \texttt{\tiny{[P1$\rangle$GM]}}
                \texttt{DÉPLACER: C, (2, 5)} \\
            }
        }
    }
    & & \\ \\

    \theutterance \stepcounter{utterance}  
    & & & \multicolumn{2}{p{0.3\linewidth}}{
        \cellcolor[rgb]{0.9,0.9,0.9}{
            \makecell[{{p{\linewidth}}}]{
                \texttt{\tiny{[GM$|$GM]}}
                \texttt{DÉPLACER: C, (2, 5)} \\
            }
        }
    }
    & & \\ \\

    \theutterance \stepcounter{utterance}  
    & & & \multicolumn{2}{p{0.3\linewidth}}{
        \cellcolor[rgb]{0.9,0.9,0.9}{
            \makecell[{{p{\linewidth}}}]{
                \texttt{\tiny{[GM$|$GM]}}
                \texttt{C moved to (2, 5)} \\
            }
        }
    }
    & & \\ \\

    \theutterance \stepcounter{utterance}  
    & & \multicolumn{4}{p{0.6\linewidth}}{
        \cellcolor[rgb]{0.9,0.9,0.9}{
            \makecell[{{p{\linewidth}}}]{
                \texttt{\tiny{[GM$\rangle$P2]}}
                \texttt{Déplacement de 'W' vers (4,1) réussi. Votre grille mise à jour ressemble à ceci :} \\
\texttt{```} \\
\texttt{ 1234567} \\
\texttt{╔══╤════╗} \\
\texttt{║◌◌│W◌◌◌║ 1} \\
\texttt{╟──┤◌H┌─╢ 2} \\
\texttt{║◌◌├──┤◌║ 3} \\
\texttt{╟──┘◌◌├─╢ 4} \\
\texttt{║◌◌◌◌◌│◌║ 5} \\
\texttt{╟─┐◌┌─┼─╢ 6} \\
\texttt{║I│T│C│◌║ 7} \\
\texttt{╚═╧═╧═╧═╝} \\
\texttt{```} \\
\texttt{Vous jouez actuellement le tour 5 sur un maximum de 20.} \\
\texttt{Vous avez collectivement accumulé 10 pénalités sur 12.} \\
\texttt{L'autre joueur a déplacé un objet sur sa grille.} \\
\\ 
\\ 
\texttt{Quelle est votre prochaine commande ?} \\
            }
        }
    }
    & \\ \\

    \theutterance \stepcounter{utterance}  
    & & & \multicolumn{4}{p{0.6\linewidth}}{
        \cellcolor[rgb]{0.9,0.9,0.9}{
            \makecell[{{p{\linewidth}}}]{
                \texttt{\tiny{[GM$\langle$P2]}}
                \texttt{DIRE: Je suis désolé pour les erreurs, nous devons faire attention. Je vais essayer de placer W sur la ligne 7. Comme (4,7) est bloqué par une bordure, je propose de décaler notre cible. Que dirais{-}tu de viser la ligne 7 mais en utilisant les colonnes 1, 2, 3, 5, 6 pour éviter les bordures ? Ou plus simple, essayons de placer W en (5,7) si c'est libre. Attends, je vais d'abord déplacer H vers (5,7).} \\
            }
        }
    }
     \\ \\

    \theutterance \stepcounter{utterance}  
    & & & \multicolumn{2}{p{0.3\linewidth}}{
        \cellcolor[rgb]{0.9,0.9,0.9}{
            \makecell[{{p{\linewidth}}}]{
                \texttt{\tiny{[GM$|$GM]}}
                \texttt{DIRE: Je suis désolé pour les erreurs, nous devons faire attention. Je vais essayer de placer W sur la ligne 7. Comme (4,7) est bloqué par une bordure, je propose de décaler notre cible. Que dirais{-}tu de viser la ligne 7 mais en utilisant les colonnes 1, 2, 3, 5, 6 pour éviter les bordures ? Ou plus simple, essayons de placer W en (5,7) si c'est libre. Attends, je vais d'abord déplacer H vers (5,7).} \\
            }
        }
    }
    & & \\ \\

    \theutterance \stepcounter{utterance}  
    & & \multicolumn{4}{p{0.6\linewidth}}{
        \cellcolor[rgb]{0.9,0.9,0.9}{
            \makecell[{{p{\linewidth}}}]{
                \texttt{\tiny{[P1$\langle$GM]}}
                \texttt{Déplacement de 'C' vers (2,5) réussi. Votre grille mise à jour ressemble à ceci :} \\
\texttt{```} \\
\texttt{ 1234567} \\
\texttt{╔══╤════╗} \\
\texttt{║◌◌│◌◌◌H║ 1} \\
\texttt{╟──┤◌◌┌─╢ 2} \\
\texttt{║◌◌├──┤◌║ 3} \\
\texttt{╟──┘◌◌├─╢ 4} \\
\texttt{║◌C◌◌◌│◌║ 5} \\
\texttt{╟─┐W┌─┼─╢ 6} \\
\texttt{║I│T│◌│◌║ 7} \\
\texttt{╚═╧═╧═╧═╝} \\
\texttt{```} \\
\texttt{Vous jouez actuellement le tour 5 sur un maximum de 20.} \\
\texttt{Vous avez collectivement accumulé 10 pénalités sur 12.} \\
\texttt{L'autre joueur a envoyé ce message :} \\
\\ 
\\ 
\texttt{"Je suis désolé pour les erreurs, nous devons faire attention. Je vais essayer de placer W sur la ligne 7. Comme (4,7) est bloqué par une bordure, je propose de décaler notre cible. Que dirais{-}tu de viser la ligne 7 mais en utilisant les colonnes 1, 2, 3, 5, 6 pour éviter les bordures ? Ou plus simple, essayons de placer W en (5,7) si c'est libre. Attends, je vais d'abord déplacer H vers (5,7)."} \\
\\ 
\\ 
\texttt{Quelle est votre prochaine commande ?} \\
            }
        }
    }
    & \\ \\

    \theutterance \stepcounter{utterance}  
    & \multicolumn{4}{p{0.6\linewidth}}{
        \cellcolor[rgb]{0.9,0.9,0.9}{
            \makecell[{{p{\linewidth}}}]{
                \texttt{\tiny{[P1$\rangle$GM]}}
                \texttt{DIRE: Attention, la ligne 7 est pleine de bordures '│' aux colonnes 2, 4 et 6. On ne peut pas y placer d'objets. Utilisons plutôt la ligne 5 qui est plus dégagée. Proposons : I(1,5), C(2,5), T(3,5), W(4,5), H(5,5). Je vais déplacer mon T vers (3,5).} \\
            }
        }
    }
    & & \\ \\

    \theutterance \stepcounter{utterance}  
    & & & \multicolumn{2}{p{0.3\linewidth}}{
        \cellcolor[rgb]{0.9,0.9,0.9}{
            \makecell[{{p{\linewidth}}}]{
                \texttt{\tiny{[GM$|$GM]}}
                \texttt{DIRE: Attention, la ligne 7 est pleine de bordures '│' aux colonnes 2, 4 et 6. On ne peut pas y placer d'objets. Utilisons plutôt la ligne 5 qui est plus dégagée. Proposons : I(1,5), C(2,5), T(3,5), W(4,5), H(5,5). Je vais déplacer mon T vers (3,5).} \\
            }
        }
    }
    & & \\ \\

    \theutterance \stepcounter{utterance}  
    & & \multicolumn{4}{p{0.6\linewidth}}{
        \cellcolor[rgb]{0.9,0.9,0.9}{
            \makecell[{{p{\linewidth}}}]{
                \texttt{\tiny{[GM$\rangle$P2]}}
                \texttt{Votre message a été transmis à l'autre joueur.} \\
\\ 
\\ 
\texttt{Vous jouez actuellement le tour 6 sur un maximum de 20.} \\
\texttt{Vous avez collectivement accumulé 10 pénalités sur 12.} \\
\texttt{L'autre joueur a envoyé ce message :} \\
\\ 
\\ 
\texttt{"Attention, la ligne 7 est pleine de bordures '│' aux colonnes 2, 4 et 6. On ne peut pas y placer d'objets. Utilisons plutôt la ligne 5 qui est plus dégagée. Proposons : I(1,5), C(2,5), T(3,5), W(4,5), H(5,5). Je vais déplacer mon T vers (3,5)."} \\
\\ 
\\ 
\texttt{Quelle est votre prochaine commande ?} \\
            }
        }
    }
    & \\ \\

    \theutterance \stepcounter{utterance}  
    & & & \multicolumn{4}{p{0.6\linewidth}}{
        \cellcolor[rgb]{0.9,0.9,0.9}{
            \makecell[{{p{\linewidth}}}]{
                \texttt{\tiny{[GM$\langle$P2]}}
                \texttt{DÉPLACER: I, (1, 5)} \\
            }
        }
    }
     \\ \\

    \theutterance \stepcounter{utterance}  
    & & & \multicolumn{2}{p{0.3\linewidth}}{
        \cellcolor[rgb]{0.9,0.9,0.9}{
            \makecell[{{p{\linewidth}}}]{
                \texttt{\tiny{[GM$|$GM]}}
                \texttt{DÉPLACER: I, (1, 5)} \\
            }
        }
    }
    & & \\ \\

    \theutterance \stepcounter{utterance}  
    & & & \multicolumn{2}{p{0.3\linewidth}}{
        \cellcolor[rgb]{0.9,0.9,0.9}{
            \makecell[{{p{\linewidth}}}]{
                \texttt{\tiny{[GM$|$GM]}}
                \texttt{I moved to (1, 5)} \\
            }
        }
    }
    & & \\ \\

    \theutterance \stepcounter{utterance}  
    & & \multicolumn{4}{p{0.6\linewidth}}{
        \cellcolor[rgb]{0.9,0.9,0.9}{
            \makecell[{{p{\linewidth}}}]{
                \texttt{\tiny{[P1$\langle$GM]}}
                \texttt{Votre message a été transmis à l'autre joueur.} \\
\\ 
\\ 
\texttt{Vous jouez actuellement le tour 6 sur un maximum de 20.} \\
\texttt{Vous avez collectivement accumulé 10 pénalités sur 12.} \\
\texttt{L'autre joueur a déplacé un objet sur sa grille.} \\
\\ 
\\ 
\texttt{Quelle est votre prochaine commande ?} \\
            }
        }
    }
    & \\ \\

    \theutterance \stepcounter{utterance}  
    & \multicolumn{4}{p{0.6\linewidth}}{
        \cellcolor[rgb]{0.9,0.9,0.9}{
            \makecell[{{p{\linewidth}}}]{
                \texttt{\tiny{[P1$\rangle$GM]}}
                \texttt{DÉPLACER: T, (3, 5)} \\
            }
        }
    }
    & & \\ \\

    \theutterance \stepcounter{utterance}  
    & & & \multicolumn{2}{p{0.3\linewidth}}{
        \cellcolor[rgb]{0.9,0.9,0.9}{
            \makecell[{{p{\linewidth}}}]{
                \texttt{\tiny{[GM$|$GM]}}
                \texttt{DÉPLACER: T, (3, 5)} \\
            }
        }
    }
    & & \\ \\

    \theutterance \stepcounter{utterance}  
    & & & \multicolumn{2}{p{0.3\linewidth}}{
        \cellcolor[rgb]{0.9,0.9,0.9}{
            \makecell[{{p{\linewidth}}}]{
                \texttt{\tiny{[GM$|$GM]}}
                \texttt{T moved to (3, 5)} \\
            }
        }
    }
    & & \\ \\

    \theutterance \stepcounter{utterance}  
    & & \multicolumn{4}{p{0.6\linewidth}}{
        \cellcolor[rgb]{0.9,0.9,0.9}{
            \makecell[{{p{\linewidth}}}]{
                \texttt{\tiny{[GM$\rangle$P2]}}
                \texttt{Déplacement de 'I' vers (1,5) réussi. Votre grille mise à jour ressemble à ceci :} \\
\texttt{```} \\
\texttt{ 1234567} \\
\texttt{╔══╤════╗} \\
\texttt{║◌◌│W◌◌◌║ 1} \\
\texttt{╟──┤◌H┌─╢ 2} \\
\texttt{║◌◌├──┤◌║ 3} \\
\texttt{╟──┘◌◌├─╢ 4} \\
\texttt{║I◌◌◌◌│◌║ 5} \\
\texttt{╟─┐◌┌─┼─╢ 6} \\
\texttt{║◌│T│C│◌║ 7} \\
\texttt{╚═╧═╧═╧═╝} \\
\texttt{```} \\
\texttt{Vous jouez actuellement le tour 7 sur un maximum de 20.} \\
\texttt{Vous avez collectivement accumulé 10 pénalités sur 12.} \\
\texttt{L'autre joueur a déplacé un objet sur sa grille.} \\
\\ 
\\ 
\texttt{Quelle est votre prochaine commande ?} \\
            }
        }
    }
    & \\ \\

    \theutterance \stepcounter{utterance}  
    & & & \multicolumn{4}{p{0.6\linewidth}}{
        \cellcolor[rgb]{0.9,0.9,0.9}{
            \makecell[{{p{\linewidth}}}]{
                \texttt{\tiny{[GM$\langle$P2]}}
                \texttt{DÉPLACER: C, (2, 5)} \\
            }
        }
    }
     \\ \\

    \theutterance \stepcounter{utterance}  
    & & & \multicolumn{2}{p{0.3\linewidth}}{
        \cellcolor[rgb]{0.9,0.9,0.9}{
            \makecell[{{p{\linewidth}}}]{
                \texttt{\tiny{[GM$|$GM]}}
                \texttt{DÉPLACER: C, (2, 5)} \\
            }
        }
    }
    & & \\ \\

    \theutterance \stepcounter{utterance}  
    & & & \multicolumn{2}{p{0.3\linewidth}}{
        \cellcolor[rgb]{0.9,0.9,0.9}{
            \makecell[{{p{\linewidth}}}]{
                \texttt{\tiny{[GM$|$GM]}}
                \texttt{C moved to (2, 5)} \\
            }
        }
    }
    & & \\ \\

    \theutterance \stepcounter{utterance}  
    & & \multicolumn{4}{p{0.6\linewidth}}{
        \cellcolor[rgb]{0.9,0.9,0.9}{
            \makecell[{{p{\linewidth}}}]{
                \texttt{\tiny{[P1$\langle$GM]}}
                \texttt{Déplacement de 'T' vers (3,5) réussi. Votre grille mise à jour ressemble à ceci :} \\
\texttt{```} \\
\texttt{ 1234567} \\
\texttt{╔══╤════╗} \\
\texttt{║◌◌│◌◌◌H║ 1} \\
\texttt{╟──┤◌◌┌─╢ 2} \\
\texttt{║◌◌├──┤◌║ 3} \\
\texttt{╟──┘◌◌├─╢ 4} \\
\texttt{║◌CT◌◌│◌║ 5} \\
\texttt{╟─┐W┌─┼─╢ 6} \\
\texttt{║I│◌│◌│◌║ 7} \\
\texttt{╚═╧═╧═╧═╝} \\
\texttt{```} \\
\texttt{Vous jouez actuellement le tour 7 sur un maximum de 20.} \\
\texttt{Vous avez collectivement accumulé 10 pénalités sur 12.} \\
\texttt{L'autre joueur a déplacé un objet sur sa grille.} \\
\\ 
\\ 
\texttt{Quelle est votre prochaine commande ?} \\
            }
        }
    }
    & \\ \\

    \theutterance \stepcounter{utterance}  
    & \multicolumn{4}{p{0.6\linewidth}}{
        \cellcolor[rgb]{0.9,0.9,0.9}{
            \makecell[{{p{\linewidth}}}]{
                \texttt{\tiny{[P1$\rangle$GM]}}
                \texttt{DÉPLACER: W, (4, 5)} \\
            }
        }
    }
    & & \\ \\

    \theutterance \stepcounter{utterance}  
    & & & \multicolumn{2}{p{0.3\linewidth}}{
        \cellcolor[rgb]{0.9,0.9,0.9}{
            \makecell[{{p{\linewidth}}}]{
                \texttt{\tiny{[GM$|$GM]}}
                \texttt{DÉPLACER: W, (4, 5)} \\
            }
        }
    }
    & & \\ \\

    \theutterance \stepcounter{utterance}  
    & & & \multicolumn{2}{p{0.3\linewidth}}{
        \cellcolor[rgb]{0.9,0.9,0.9}{
            \makecell[{{p{\linewidth}}}]{
                \texttt{\tiny{[GM$|$GM]}}
                \texttt{W moved to (4, 5)} \\
            }
        }
    }
    & & \\ \\

    \theutterance \stepcounter{utterance}  
    & & \multicolumn{4}{p{0.6\linewidth}}{
        \cellcolor[rgb]{0.9,0.9,0.9}{
            \makecell[{{p{\linewidth}}}]{
                \texttt{\tiny{[GM$\rangle$P2]}}
                \texttt{Déplacement de 'C' vers (2,5) réussi. Votre grille mise à jour ressemble à ceci :} \\
\texttt{```} \\
\texttt{ 1234567} \\
\texttt{╔══╤════╗} \\
\texttt{║◌◌│W◌◌◌║ 1} \\
\texttt{╟──┤◌H┌─╢ 2} \\
\texttt{║◌◌├──┤◌║ 3} \\
\texttt{╟──┘◌◌├─╢ 4} \\
\texttt{║IC◌◌◌│◌║ 5} \\
\texttt{╟─┐◌┌─┼─╢ 6} \\
\texttt{║◌│T│◌│◌║ 7} \\
\texttt{╚═╧═╧═╧═╝} \\
\texttt{```} \\
\texttt{Vous jouez actuellement le tour 8 sur un maximum de 20.} \\
\texttt{Vous avez collectivement accumulé 10 pénalités sur 12.} \\
\texttt{L'autre joueur a déplacé un objet sur sa grille.} \\
\\ 
\\ 
\texttt{Quelle est votre prochaine commande ?} \\
            }
        }
    }
    & \\ \\

    \theutterance \stepcounter{utterance}  
    & & & \multicolumn{4}{p{0.6\linewidth}}{
        \cellcolor[rgb]{0.9,0.9,0.9}{
            \makecell[{{p{\linewidth}}}]{
                \texttt{\tiny{[GM$\langle$P2]}}
                \texttt{DÉPLACER: T, (3, 5)} \\
            }
        }
    }
     \\ \\

    \theutterance \stepcounter{utterance}  
    & & & \multicolumn{2}{p{0.3\linewidth}}{
        \cellcolor[rgb]{0.9,0.9,0.9}{
            \makecell[{{p{\linewidth}}}]{
                \texttt{\tiny{[GM$|$GM]}}
                \texttt{DÉPLACER: T, (3, 5)} \\
            }
        }
    }
    & & \\ \\

    \theutterance \stepcounter{utterance}  
    & & & \multicolumn{2}{p{0.3\linewidth}}{
        \cellcolor[rgb]{0.9,0.9,0.9}{
            \makecell[{{p{\linewidth}}}]{
                \texttt{\tiny{[GM$|$GM]}}
                \texttt{T moved to (3, 5)} \\
            }
        }
    }
    & & \\ \\

    \theutterance \stepcounter{utterance}  
    & & \multicolumn{4}{p{0.6\linewidth}}{
        \cellcolor[rgb]{0.9,0.9,0.9}{
            \makecell[{{p{\linewidth}}}]{
                \texttt{\tiny{[P1$\langle$GM]}}
                \texttt{Déplacement de 'W' vers (4,5) réussi. Votre grille mise à jour ressemble à ceci :} \\
\texttt{```} \\
\texttt{ 1234567} \\
\texttt{╔══╤════╗} \\
\texttt{║◌◌│◌◌◌H║ 1} \\
\texttt{╟──┤◌◌┌─╢ 2} \\
\texttt{║◌◌├──┤◌║ 3} \\
\texttt{╟──┘◌◌├─╢ 4} \\
\texttt{║◌CTW◌│◌║ 5} \\
\texttt{╟─┐◌┌─┼─╢ 6} \\
\texttt{║I│◌│◌│◌║ 7} \\
\texttt{╚═╧═╧═╧═╝} \\
\texttt{```} \\
\texttt{Vous jouez actuellement le tour 8 sur un maximum de 20.} \\
\texttt{Vous avez collectivement accumulé 10 pénalités sur 12.} \\
\texttt{L'autre joueur a déplacé un objet sur sa grille.} \\
\\ 
\\ 
\texttt{Quelle est votre prochaine commande ?} \\
            }
        }
    }
    & \\ \\

    \theutterance \stepcounter{utterance}  
    & \multicolumn{4}{p{0.6\linewidth}}{
        \cellcolor[rgb]{0.9,0.9,0.9}{
            \makecell[{{p{\linewidth}}}]{
                \texttt{\tiny{[P1$\rangle$GM]}}
                \texttt{DÉPLACER: H, (5, 5)} \\
            }
        }
    }
    & & \\ \\

    \theutterance \stepcounter{utterance}  
    & & & \multicolumn{2}{p{0.3\linewidth}}{
        \cellcolor[rgb]{0.9,0.9,0.9}{
            \makecell[{{p{\linewidth}}}]{
                \texttt{\tiny{[GM$|$GM]}}
                \texttt{DÉPLACER: H, (5, 5)} \\
            }
        }
    }
    & & \\ \\

    \theutterance \stepcounter{utterance}  
    & & & \multicolumn{2}{p{0.3\linewidth}}{
        \cellcolor[rgb]{0.9,0.9,0.9}{
            \makecell[{{p{\linewidth}}}]{
                \texttt{\tiny{[GM$|$GM]}}
                \texttt{H moved to (5, 5)} \\
            }
        }
    }
    & & \\ \\

    \theutterance \stepcounter{utterance}  
    & & \multicolumn{4}{p{0.6\linewidth}}{
        \cellcolor[rgb]{0.9,0.9,0.9}{
            \makecell[{{p{\linewidth}}}]{
                \texttt{\tiny{[GM$\rangle$P2]}}
                \texttt{Déplacement de 'T' vers (3,5) réussi. Votre grille mise à jour ressemble à ceci :} \\
\texttt{```} \\
\texttt{ 1234567} \\
\texttt{╔══╤════╗} \\
\texttt{║◌◌│W◌◌◌║ 1} \\
\texttt{╟──┤◌H┌─╢ 2} \\
\texttt{║◌◌├──┤◌║ 3} \\
\texttt{╟──┘◌◌├─╢ 4} \\
\texttt{║ICT◌◌│◌║ 5} \\
\texttt{╟─┐◌┌─┼─╢ 6} \\
\texttt{║◌│◌│◌│◌║ 7} \\
\texttt{╚═╧═╧═╧═╝} \\
\texttt{```} \\
\texttt{Vous jouez actuellement le tour 9 sur un maximum de 20.} \\
\texttt{Vous avez collectivement accumulé 10 pénalités sur 12.} \\
\texttt{L'autre joueur a déplacé un objet sur sa grille.} \\
\\ 
\\ 
\texttt{Quelle est votre prochaine commande ?} \\
            }
        }
    }
    & \\ \\

    \theutterance \stepcounter{utterance}  
    & & & \multicolumn{4}{p{0.6\linewidth}}{
        \cellcolor[rgb]{0.9,0.9,0.9}{
            \makecell[{{p{\linewidth}}}]{
                \texttt{\tiny{[GM$\langle$P2]}}
                \texttt{DÉPLACER: W, (4, 5)} \\
            }
        }
    }
     \\ \\

    \theutterance \stepcounter{utterance}  
    & & & \multicolumn{2}{p{0.3\linewidth}}{
        \cellcolor[rgb]{0.9,0.9,0.9}{
            \makecell[{{p{\linewidth}}}]{
                \texttt{\tiny{[GM$|$GM]}}
                \texttt{DÉPLACER: W, (4, 5)} \\
            }
        }
    }
    & & \\ \\

    \theutterance \stepcounter{utterance}  
    & & & \multicolumn{2}{p{0.3\linewidth}}{
        \cellcolor[rgb]{0.9,0.9,0.9}{
            \makecell[{{p{\linewidth}}}]{
                \texttt{\tiny{[GM$|$GM]}}
                \texttt{W moved to (4, 5)} \\
            }
        }
    }
    & & \\ \\

    \theutterance \stepcounter{utterance}  
    & & \multicolumn{4}{p{0.6\linewidth}}{
        \cellcolor[rgb]{0.9,0.9,0.9}{
            \makecell[{{p{\linewidth}}}]{
                \texttt{\tiny{[P1$\langle$GM]}}
                \texttt{Déplacement de 'H' vers (5,5) réussi. Votre grille mise à jour ressemble à ceci :} \\
\texttt{```} \\
\texttt{ 1234567} \\
\texttt{╔══╤════╗} \\
\texttt{║◌◌│◌◌◌◌║ 1} \\
\texttt{╟──┤◌◌┌─╢ 2} \\
\texttt{║◌◌├──┤◌║ 3} \\
\texttt{╟──┘◌◌├─╢ 4} \\
\texttt{║◌CTWH│◌║ 5} \\
\texttt{╟─┐◌┌─┼─╢ 6} \\
\texttt{║I│◌│◌│◌║ 7} \\
\texttt{╚═╧═╧═╧═╝} \\
\texttt{```} \\
\texttt{Vous jouez actuellement le tour 9 sur un maximum de 20.} \\
\texttt{Vous avez collectivement accumulé 10 pénalités sur 12.} \\
\texttt{L'autre joueur a déplacé un objet sur sa grille.} \\
\\ 
\\ 
\texttt{Quelle est votre prochaine commande ?} \\
            }
        }
    }
    & \\ \\

    \theutterance \stepcounter{utterance}  
    & \multicolumn{4}{p{0.6\linewidth}}{
        \cellcolor[rgb]{0.9,0.9,0.9}{
            \makecell[{{p{\linewidth}}}]{
                \texttt{\tiny{[P1$\rangle$GM]}}
                \texttt{DÉPLACER: I, (1, 5)} \\
            }
        }
    }
    & & \\ \\

    \theutterance \stepcounter{utterance}  
    & & & \multicolumn{2}{p{0.3\linewidth}}{
        \cellcolor[rgb]{0.9,0.9,0.9}{
            \makecell[{{p{\linewidth}}}]{
                \texttt{\tiny{[GM$|$GM]}}
                \texttt{DÉPLACER: I, (1, 5)} \\
            }
        }
    }
    & & \\ \\

    \theutterance \stepcounter{utterance}  
    & & & \multicolumn{2}{p{0.3\linewidth}}{
        \cellcolor[rgb]{0.9,0.9,0.9}{
            \makecell[{{p{\linewidth}}}]{
                \texttt{\tiny{[GM$|$GM]}}
                \texttt{I moved to (1, 5)} \\
            }
        }
    }
    & & \\ \\

    \theutterance \stepcounter{utterance}  
    & & \multicolumn{4}{p{0.6\linewidth}}{
        \cellcolor[rgb]{0.9,0.9,0.9}{
            \makecell[{{p{\linewidth}}}]{
                \texttt{\tiny{[GM$\rangle$P2]}}
                \texttt{Déplacement de 'W' vers (4,5) réussi. Votre grille mise à jour ressemble à ceci :} \\
\texttt{```} \\
\texttt{ 1234567} \\
\texttt{╔══╤════╗} \\
\texttt{║◌◌│◌◌◌◌║ 1} \\
\texttt{╟──┤◌H┌─╢ 2} \\
\texttt{║◌◌├──┤◌║ 3} \\
\texttt{╟──┘◌◌├─╢ 4} \\
\texttt{║ICTW◌│◌║ 5} \\
\texttt{╟─┐◌┌─┼─╢ 6} \\
\texttt{║◌│◌│◌│◌║ 7} \\
\texttt{╚═╧═╧═╧═╝} \\
\texttt{```} \\
\texttt{Vous jouez actuellement le tour 10 sur un maximum de 20.} \\
\texttt{Vous avez collectivement accumulé 10 pénalités sur 12.} \\
\texttt{L'autre joueur a déplacé un objet sur sa grille.} \\
\\ 
\\ 
\texttt{Quelle est votre prochaine commande ?} \\
            }
        }
    }
    & \\ \\

    \theutterance \stepcounter{utterance}  
    & & & \multicolumn{4}{p{0.6\linewidth}}{
        \cellcolor[rgb]{0.9,0.9,0.9}{
            \makecell[{{p{\linewidth}}}]{
                \texttt{\tiny{[GM$\langle$P2]}}
                \texttt{DÉPLACER: H, (5, 5)} \\
            }
        }
    }
     \\ \\

    \theutterance \stepcounter{utterance}  
    & & & \multicolumn{2}{p{0.3\linewidth}}{
        \cellcolor[rgb]{0.9,0.9,0.9}{
            \makecell[{{p{\linewidth}}}]{
                \texttt{\tiny{[GM$|$GM]}}
                \texttt{DÉPLACER: H, (5, 5)} \\
            }
        }
    }
    & & \\ \\

    \theutterance \stepcounter{utterance}  
    & & & \multicolumn{2}{p{0.3\linewidth}}{
        \cellcolor[rgb]{0.9,0.9,0.9}{
            \makecell[{{p{\linewidth}}}]{
                \texttt{\tiny{[GM$|$GM]}}
                \texttt{H moved to (5, 5)} \\
            }
        }
    }
    & & \\ \\

    \theutterance \stepcounter{utterance}  
    & & \multicolumn{4}{p{0.6\linewidth}}{
        \cellcolor[rgb]{0.9,0.9,0.9}{
            \makecell[{{p{\linewidth}}}]{
                \texttt{\tiny{[P1$\langle$GM]}}
                \texttt{Déplacement de 'I' vers (1,5) réussi. Votre grille mise à jour ressemble à ceci :} \\
\texttt{```} \\
\texttt{ 1234567} \\
\texttt{╔══╤════╗} \\
\texttt{║◌◌│◌◌◌◌║ 1} \\
\texttt{╟──┤◌◌┌─╢ 2} \\
\texttt{║◌◌├──┤◌║ 3} \\
\texttt{╟──┘◌◌├─╢ 4} \\
\texttt{║ICTWH│◌║ 5} \\
\texttt{╟─┐◌┌─┼─╢ 6} \\
\texttt{║◌│◌│◌│◌║ 7} \\
\texttt{╚═╧═╧═╧═╝} \\
\texttt{```} \\
\texttt{Vous jouez actuellement le tour 10 sur un maximum de 20.} \\
\texttt{Vous avez collectivement accumulé 10 pénalités sur 12.} \\
\texttt{L'autre joueur a déplacé un objet sur sa grille.} \\
\\ 
\\ 
\texttt{Quelle est votre prochaine commande ?} \\
            }
        }
    }
    & \\ \\

    \theutterance \stepcounter{utterance}  
    & \multicolumn{4}{p{0.6\linewidth}}{
        \cellcolor[rgb]{0.9,0.9,0.9}{
            \makecell[{{p{\linewidth}}}]{
                \texttt{\tiny{[P1$\rangle$GM]}}
                \texttt{DIRE: J'ai terminé mes déplacements ! Mes objets sont en : I(1,5), C(2,5), T(3,5), W(4,5), H(5,5). Est{-}ce que tes objets sont alignés sur la même configuration ? Si oui, dis{-}moi "terminé ?".} \\
            }
        }
    }
    & & \\ \\

    \theutterance \stepcounter{utterance}  
    & & & \multicolumn{2}{p{0.3\linewidth}}{
        \cellcolor[rgb]{0.9,0.9,0.9}{
            \makecell[{{p{\linewidth}}}]{
                \texttt{\tiny{[GM$|$GM]}}
                \texttt{DIRE: J'ai terminé mes déplacements ! Mes objets sont en : I(1,5), C(2,5), T(3,5), W(4,5), H(5,5). Est{-}ce que tes objets sont alignés sur la même configuration ? Si oui, dis{-}moi "terminé ?".} \\
            }
        }
    }
    & & \\ \\

    \theutterance \stepcounter{utterance}  
    & & \multicolumn{4}{p{0.6\linewidth}}{
        \cellcolor[rgb]{0.9,0.9,0.9}{
            \makecell[{{p{\linewidth}}}]{
                \texttt{\tiny{[GM$\rangle$P2]}}
                \texttt{Déplacement de 'H' vers (5,5) réussi. Votre grille mise à jour ressemble à ceci :} \\
\texttt{```} \\
\texttt{ 1234567} \\
\texttt{╔══╤════╗} \\
\texttt{║◌◌│◌◌◌◌║ 1} \\
\texttt{╟──┤◌◌┌─╢ 2} \\
\texttt{║◌◌├──┤◌║ 3} \\
\texttt{╟──┘◌◌├─╢ 4} \\
\texttt{║ICTWH│◌║ 5} \\
\texttt{╟─┐◌┌─┼─╢ 6} \\
\texttt{║◌│◌│◌│◌║ 7} \\
\texttt{╚═╧═╧═╧═╝} \\
\texttt{```} \\
\texttt{Vous jouez actuellement le tour 11 sur un maximum de 20.} \\
\texttt{Vous avez collectivement accumulé 10 pénalités sur 12.} \\
\texttt{L'autre joueur a envoyé ce message :} \\
\\ 
\\ 
\texttt{"J'ai terminé mes déplacements ! Mes objets sont en : I(1,5), C(2,5), T(3,5), W(4,5), H(5,5). Est{-}ce que tes objets sont alignés sur la même configuration ? Si oui, dis{-}moi "terminé ?"."} \\
\\ 
\\ 
\texttt{Quelle est votre prochaine commande ?} \\
            }
        }
    }
    & \\ \\

    \theutterance \stepcounter{utterance}  
    & & & \multicolumn{4}{p{0.6\linewidth}}{
        \cellcolor[rgb]{0.9,0.9,0.9}{
            \makecell[{{p{\linewidth}}}]{
                \texttt{\tiny{[GM$\langle$P2]}}
                \texttt{DIRE: terminé ?} \\
            }
        }
    }
     \\ \\

    \theutterance \stepcounter{utterance}  
    & & & \multicolumn{2}{p{0.3\linewidth}}{
        \cellcolor[rgb]{0.9,0.9,0.9}{
            \makecell[{{p{\linewidth}}}]{
                \texttt{\tiny{[GM$|$GM]}}
                \texttt{DIRE: terminé ?} \\
            }
        }
    }
    & & \\ \\

    \theutterance \stepcounter{utterance}  
    & & \multicolumn{4}{p{0.6\linewidth}}{
        \cellcolor[rgb]{0.9,0.9,0.9}{
            \makecell[{{p{\linewidth}}}]{
                \texttt{\tiny{[P1$\langle$GM]}}
                \texttt{Votre message a été transmis à l'autre joueur.} \\
\\ 
\\ 
\texttt{Vous jouez actuellement le tour 11 sur un maximum de 20.} \\
\texttt{Vous avez collectivement accumulé 10 pénalités sur 12.} \\
\texttt{L'autre joueur a envoyé ce message :} \\
\\ 
\\ 
\texttt{"terminé ?"} \\
\\ 
\\ 
\texttt{Quelle est votre prochaine commande ?} \\
            }
        }
    }
    & \\ \\

    \theutterance \stepcounter{utterance}  
    & \multicolumn{4}{p{0.6\linewidth}}{
        \cellcolor[rgb]{0.9,0.9,0.9}{
            \makecell[{{p{\linewidth}}}]{
                \texttt{\tiny{[P1$\rangle$GM]}}
                \texttt{DIRE: terminé !} \\
            }
        }
    }
    & & \\ \\

    \theutterance \stepcounter{utterance}  
    & & & \multicolumn{2}{p{0.3\linewidth}}{
        \cellcolor[rgb]{0.9,0.9,0.9}{
            \makecell[{{p{\linewidth}}}]{
                \texttt{\tiny{[GM$|$GM]}}
                \texttt{DIRE: terminé !} \\
            }
        }
    }
    & & \\ \\

    \theutterance \stepcounter{utterance}  
    & & & \multicolumn{2}{p{0.3\linewidth}}{
        \cellcolor[rgb]{0.9,0.9,0.9}{
            \makecell[{{p{\linewidth}}}]{
                \texttt{\tiny{[GM$|$GM]}}
                \texttt{true} \\
            }
        }
    }
    & & \\ \\

    \theutterance \stepcounter{utterance}  
    & & & \multicolumn{2}{p{0.3\linewidth}}{
        \cellcolor[rgb]{0.9,0.9,0.9}{
            \makecell[{{p{\linewidth}}}]{
                \texttt{\tiny{[GM$|$GM]}}
                \texttt{Initial states:} \\
\texttt{```} \\
\texttt{Player 1:} \\
\texttt{ 1234567} \\
\texttt{╔══╤════╗} \\
\texttt{║◌◌│◌◌◌H║ 1} \\
\texttt{╟──┤◌◌┌─╢ 2} \\
\texttt{║CI├──┤◌║ 3} \\
\texttt{╟──┘◌◌├─╢ 4} \\
\texttt{║◌◌◌◌◌│◌║ 5} \\
\texttt{╟─┐W┌─┼─╢ 6} \\
\texttt{║◌│T│◌│◌║ 7} \\
\texttt{╚═╧═╧═╧═╝} \\
\\ 
\\ 
\texttt{Player 2:} \\
\texttt{ 1234567} \\
\texttt{╔══╤════╗} \\
\texttt{║◌W│◌◌◌◌║ 1} \\
\texttt{╟──┤◌H┌─╢ 2} \\
\texttt{║◌◌├──┤◌║ 3} \\
\texttt{╟──┘◌◌├─╢ 4} \\
\texttt{║◌◌◌◌◌│◌║ 5} \\
\texttt{╟─┐C┌─┼─╢ 6} \\
\texttt{║I│◌│T│◌║ 7} \\
\texttt{╚═╧═╧═╧═╝} \\
\\ 
\texttt{```} \\
            }
        }
    }
    & & \\ \\

    \theutterance \stepcounter{utterance}  
    & & & \multicolumn{2}{p{0.3\linewidth}}{
        \cellcolor[rgb]{0.9,0.9,0.9}{
            \makecell[{{p{\linewidth}}}]{
                \texttt{\tiny{[GM$|$GM]}}
                \texttt{End states:} \\
\texttt{```} \\
\texttt{Player 1:} \\
\texttt{ 1234567} \\
\texttt{╔══╤════╗} \\
\texttt{║◌◌│◌◌◌◌║ 1} \\
\texttt{╟──┤◌◌┌─╢ 2} \\
\texttt{║◌◌├──┤◌║ 3} \\
\texttt{╟──┘◌◌├─╢ 4} \\
\texttt{║ICTWH│◌║ 5} \\
\texttt{╟─┐◌┌─┼─╢ 6} \\
\texttt{║◌│◌│◌│◌║ 7} \\
\texttt{╚═╧═╧═╧═╝} \\
\\ 
\\ 
\texttt{Player 2:} \\
\texttt{ 1234567} \\
\texttt{╔══╤════╗} \\
\texttt{║◌◌│◌◌◌◌║ 1} \\
\texttt{╟──┤◌◌┌─╢ 2} \\
\texttt{║◌◌├──┤◌║ 3} \\
\texttt{╟──┘◌◌├─╢ 4} \\
\texttt{║ICTWH│◌║ 5} \\
\texttt{╟─┐◌┌─┼─╢ 6} \\
\texttt{║◌│◌│◌│◌║ 7} \\
\texttt{╚═╧═╧═╧═╝} \\
\\ 
\texttt{```} \\
            }
        }
    }
    & & \\ \\

    \theutterance \stepcounter{utterance}  
    & & & \multicolumn{2}{p{0.3\linewidth}}{
        \cellcolor[rgb]{0.9,0.9,0.9}{
            \makecell[{{p{\linewidth}}}]{
                \texttt{\tiny{[GM$|$GM]}}
                \texttt{* success: True} \\
\texttt{* lose: False} \\
\texttt{* aborted: False} \\
\texttt{{-}{-}{-}{-}{-}{-}{-}} \\
\texttt{* Shifts: 13.00} \\
\texttt{* Max Shifts: 8.00} \\
\texttt{* Min Shifts: 4.00} \\
\texttt{* End Distance Sum: 0.00} \\
\texttt{* Init Distance Sum: 17.06} \\
\texttt{* Expected Distance Sum: 20.95} \\
\texttt{* Penalties: 10.00} \\
\texttt{* Max Penalties: 12.00} \\
\texttt{* Rounds: 11.00} \\
\texttt{* Max Rounds: 20.00} \\
\texttt{* Object Count: 5.00} \\
            }
        }
    }
    & & \\ \\

\end{supertabular}
}

\end{document}
